\chapter{The Polyxan Compiler Correctness Framework}
\label{ch:38}

In the previous chapter we introduced verified dynamics: the notion that any admissible transformation of a coupled configuration
[
(\Phi,\mathbf{v},S;,\mathcal{G},\iota)
]
must satisfy the modal predicates
[
\Box_{\mathrm{RSVP}} \qquad\text{and}\qquad \Box_{\mathrm{Polyx}}.
]
These modalities express, respectively, harmonic inevitability and quine-stability.

We now develop the full compiler correctness framework for Polyxan rewriting.
This framework provides the formal bridge between:

\begin{enumerate}
\item the mathematical EHR rewriting system, defined in Part~II; and
\item the executable Polyxan compiler, which produces verified rewrite sequences in the implementation layer of Part~V.
\end{enumerate}

The central goal of the chapter is to prove that compilation and mathematical rewriting coincide.
All correctness results follow from the fundamental curvature--entropy structure of the Polyxan system and the modal verification layer of Part~IV.

\section{Polyxan Source Graphs and the Rewriting Alphabet}

We begin by recalling the structure of a Polyxan hypergraph.

\begin{definition}[Typed Polyxan Hypergraph]
A typed Polyxan hypergraph is a quintuple
[
\mathcal{G}=(V,E,\tau,\mathrm{tag},\kappa_{\mathcal{G}})
]
where:
\begin{itemize}
\item $V$ are vertices,
\item $E$ are hyperedges with incidence relations,
\item $\tau : E\to \mathcal{T}$ assigns types from a finite type alphabet,
\item $\mathrm{tag} : E\to \mathcal{H}*{\mathrm{tag}}$ assigns entropy tags,
\item $\kappa*{\mathcal{G}} : E\to \mathbb{Z}_{\ge 0}$ assigns primitive Polyxan curvature.
\end{itemize}
\end{definition}

The EHR rewriting alphabet consists of three verified operations:
[
\mathrm{E} \ (\text{Extraction}), \qquad
\mathrm{H} \ (\text{Hoisting}), \qquad
\mathrm{R} \ (\text{Curvature Reduction}).
]

Each rule reduces the Polyxan curvature $\mathcal{H}_0[\mathcal{G}]$ (Chapter~25).
Thus, the system enjoys strong normalization.

\section{Compiler Semantics}

The compiler must translate a Polyxan hypergraph into an executable rewrite schedule.

\begin{definition}[Polyxan Compiler]
A Polyxan compiler is a total function
[
\mathsf{compile} : \mathbf{Polyx} \to \Lambda_{\mathrm{EHR}}
]
where $\Lambda_{\mathrm{EHR}}$ is the set of finite sequences of EHR rewrite steps, such that:
\begin{enumerate}
\item each step reduces curvature or preserves it only at a quine-bound boundary;
\item the final graph is the quine-stable normal form;
\item the sequence is provably quine-stable under $\Box_{\mathrm{Polyx}}$.
\end{enumerate}
\end{definition}

The compiler is not allowed to invent new rules; it must produce only verified EHR steps.

\begin{proposition}[Type Preservation]
$\mathsf{compile}$ preserves vertex and edge types, entropy tags, and incidence structure up to verified rewriting:
[
\mathcal{G};\leadsto_{\mathsf{compile}};\mathcal{G}'
\quad\Longrightarrow\quad
\mathcal{G}'\simeq \mathcal{G} \text{ modulo verified equivalence.}
]
\end{proposition}

\begin{proof}
All EHR rules preserve types and tags by construction (Chapter~22).
Since the compiler emits only EHR steps, type preservation follows immediately.
\end{proof}

\section{The Curvature Descent Invariant}

Central to correctness is the monotonicity of curvature.

\begin{lemma}[Curvature Descent]
For all Polyxan hypergraphs $\mathcal{G}$,
[
\mathcal{H}*0[\mathcal{G}*{i+1}] \le \mathcal{H}*0[\mathcal{G}*{i}]
]
for each rewrite step in $\mathsf{compile}(\mathcal{G})$, with equality only at quine boundaries.
\end{lemma}

\begin{proof}
Each EHR rule reduces curvature strictly unless the local region already satisfies quine-stability constraints.
This follows from the curvature-reduction theorem (Chapter~25).
\end{proof}

\begin{corollary}[Strong Normalisation]
Every compiled sequence is finite.
\end{corollary}

\section{Compiler Correctness}

We now prove the main correctness statement.

\begin{theorem}[Compiler Correctness]
\label{thm:compiler-correctness}
For every typed Polyxan hypergraph $\mathcal{G}$,
[
\mathsf{compile}(\mathcal{G})
]
terminates in a finite number of steps, producing a unique normal form that is identical to the normal form produced by the canonical EHR rewriting system:
[
\mathrm{NF}_{\mathsf{compile}}(\mathcal{G})
===========================================

\mathrm{NF}_{\mathrm{EHR}}(\mathcal{G}).
]
\end{theorem}

\begin{proof}
There are three properties to verify.

\textbf{1. Termination.}
By curvature descent, every compiler-generated rewrite chain is finite.

\textbf{2. Confluence.}
EHR rewriting is confluent (Chapter~24).
Thus, any terminating rewrite strategy yields the same normal form.

\textbf{3. Faithfulness.}
The compiler emits only rules in the EHR alphabet, hence it operates within the same rewriting theory.

Since both systems reduce curvature strictly (except at quine fixed points), their reduction sequences terminate at the same minimal-curvature configuration---the unique quine-stable normal form.
\end{proof}

\section{Modal Correctness of Compilation}

Because Part~IV introduced modal correctness layers, we must show that compilation is modal-safe.

\begin{theorem}[Modal Compiler Correctness]
[
\Box_{\mathrm{Polyx}}\bigl(\mathsf{compile}(\mathcal{G})\bigr)
]
for every Polyxan hypergraph $\mathcal{G}$.
\end{theorem}

\begin{proof}
Each step emitted by the compiler is an admissible EHR rewrite step.
By definition of $\Box_{\mathrm{Polyx}}$, any property preserved under all EHR steps is true after compilation.
Thus the entire compiled sequence is modal-safe.
\end{proof}

Combining this with stratified soundness and the harmonic embedding layer (Chapters~31--32), we obtain:

\begin{corollary}[End-to-End Verified Compilation]
If $(\Phi,\mathbf{v},S;\mathcal{G},\iota)$ is verified, then
[
(\Phi,\mathbf{v},S;,\mathrm{NF}_{\mathsf{compile}}(\mathcal{G}),\iota')
]
is also verified.
\end{corollary}

\section{Interaction with the Embedding Layer}

The compiled rewrite sequence is not purely algebraic.
It interacts with embedding flow.

\begin{theorem}[Compatibility with Embedding Flow]
Let
[
\mathcal{G} \to \mathcal{G}'
]
be a compiler-emitted EHR step.
Then the embedding flow produces a harmonic transition
[
\iota \to \iota'
]
such that
[
\mathcal{E}[\iota'] \le \mathcal{E}[\iota].
]
\end{theorem}

\begin{proof}
Every EHR step corresponds to a continuous topology-changing event in the embedding (Discrete-to-Continuous Morphogenesis, Chapter~33).
These events strictly decrease embedding energy.
Hence the flow is compatible with compilation.
\end{proof}

\section{Consequences for the Global System}

The correctness of the Polyxan compiler secures the entire implementation stack:

\begin{enumerate}
\item verified rewriting,
\item verified embedding evolution,
\item verified totality of the RSVP engine, and
\item verified convergence to the universal quine.
\end{enumerate}

\begin{corollary}[Global Verified Correctness]
Every compiled Polyxan program necessarily converges, via the embedding flow, to the universal media quine $\mathfrak{U}$.
\end{corollary}

\section{Closing Cadence}

The compiler is not merely an executable device.
It is a verified bridge between algebra, geometry, and meaning.
Because every rewrite it emits is modal-safe, curvature-decreasing, and harmonically executable, the Polyxan compiler is the formal guarantee that:

\begin{quote}
No semantic structure can escape its final normal form.
Every program reduces---provably and inevitably---to its quine-stable truth.
\end{quote}

